\section{本章小结}
\label{sec:mcerr-summary}

Monte Carlo估计$\hat{I}_N$的误差总是分成偏差与随机波动两部分。偏差只取决于各个样本的分布是否等于$\pi$，与样本是否独立无关；波动取决于样本之间的依赖与被积函数的矩。两部分必须分别评价：增加样本数只压低波动，从不消除偏差，而偏差又无法从输出本身读出。

\paragraph{各种来源的误差}
表~\ref{tab:mcerr-summary}按样本来源列出本章的主要结论。

\begin{BookInlineTable}
 \centering
 \caption{不同样本来源下的误差}
 \label{tab:mcerr-summary}
 \begin{tabular}{lll}
  \toprule
  样本来源 & 偏差 & 波动的尺度 \\
  \midrule
  独立、精确 & 零 & $\sigma_f^{2}/N$，需二阶矩 \\
  独立、近似$\tilde{\pi}$ & $\mathbb{E}_{\tilde{\pi}}f-\mathbb{E}_{\pi}f$，不随$N$变 & $\operatorname{Var}_{\tilde{\pi}}(f)/N$ \\
  链、从一般初值出发 & 由边缘分布的收敛速度控制 & 同下 \\
  链、平稳 & 零 & 渐近为$\sigma_{\mathrm{LR}}^{2}/N$ \\
  整组随机化 & 零（各点边缘正确时） & $\operatorname{Var}(Q)/R$ \\
  \bottomrule
 \end{tabular}
\end{BookInlineTable}

\paragraph{保证的类型}
独立样本有两类保证。Chebyshev界对每个$N$成立，只需二阶矩，但偏松；渐近误差区间更紧，但只在$N\to\infty$时给出名义覆盖，稀有的大贡献可以使所需的$N$大到无法从样本中察觉。只有一阶矩时大数律仍成立，但误差可以按慢于$N^{-1/2}$的速度缩小。

近似生成的偏差由总变差控制，前提是被积函数有界；无界时还要控制两种分布的尾部或矩。近似转移先要确认存在不变分布，再估计它与$\pi$的差别。

链平均的初值偏差需要边缘分布趋于$\pi$的速度估计，舍弃初始一段只有在这种估计下才有依据，对无界函数还需要一致可积。平稳链平均的方差由自协方差决定，长期方差是它的极限，有效样本量是它的换算，三者都依赖被积函数。链平均的误差区间同时需要中心极限定理、正的长期方差和一致的长期方差估计；运行诊断只能发现问题，不能证明链已遍历目标。

一次产生一组相关点的方法，以独立重复的整组为观测单位，用组间的样本方差估计误差，可靠程度由重复数决定。

\paragraph{与后续各章的关系}
第~\ref{chap:variance-reduction}章改变抽样分布、估计公式或抽样机制来减小波动，每一种改变都会改变估计量的分布，因而都要按本章的方式重新检查偏差、矩条件和样本之间的依赖。第~\ref{chap:expectation-gradient-estimation}章估计期望的导数时，近似采样带来的偏差同样不能靠样本数消除。第~\ref{chap:sampling-budget-method-selection}章在给定预算下分配样本，所比较的正是本章给出的偏差与方差。
